\scanpage{454}

formules qui ont un sens seulement si on suppose donnés des plongements (ou du moins morphismes)

\diagram{
\[
(155)\quad
\left\{
\begin{tikzcd}[ampersand replacement=\&,column sep=large]
T_\rho \arrow[r,hook] \& G \\
T_\sigma \arrow[r,hook] \& G
\end{tikzcd}
\right.
\qquad
\begin{gathered}
\text{et un sous-groupe}\\
\begin{tikzcd}[ampersand replacement=\&,column sep=large]
T_0 \arrow[r,hook] \& G
\end{tikzcd}
\end{gathered}.
\]
}

Quand on a de tels plongements, les formules (154) ont un sens, et définissent \deleted{\unsure{toutes les}} (pour les quantités (152) connues) les quantités (153) en termes de l'une \deleted{\unsure{d'elles}}\inserted{d'elles}, disons $\underline u$, ou \deleted{\illegible{1}} $f$, ou $\alpha$, ou $\beta$ (mais pas en termes de $\alpha_0$ ou de $\beta_0$).

En fait, on dispose d'un sous groupe

\diagram{
\[
\begin{array}{r@{\;}c@{\;}l}
\begin{gathered}(156)\\[2em]\deleted{\text{\unsure{tel}}}\\[3em]{}\end{gathered}
&\begin{tikzcd}[ampersand replacement=\&,row sep=3em]
{}\arrow[dd]\\{}\\{}
\end{tikzcd}
&\begin{gathered}
T_0\subset G\\
\text{tel que}\quad
\left\{\begin{array}{l}
\begin{tikzcd}[ampersand replacement=\&,column sep=large]
T_0 \arrow[r,"{\substack{\delta_G|T_0\\\sim}}"] \& G_m
\end{tikzcd}\\[5pt]
\text{i.e.}\quad G\simeq T_0.SG\\
\hspace{6em}\text{produit }1/2\text{ direct},
\end{array}\right.
\end{gathered}
\end{array}
\]
}

de sorte que \deleted{la quantité}\inserted{pour $U\in\Gamma G$,} existe un \uline{unique} \deleted{$f\in\Gamma SG$ tel que}\inserted{couple $(f,u)\in\Gamma(SG\times T_0)$ avec}
\[
U=f^{-1}\underline u.
\]
d'où des uniques $\alpha,\beta,\alpha_0,\beta_0\in\Gamma G$ satisfaisant (avec $f,u$ précédents) à (154) :
\[
(157)\quad
\alpha=\underline u\,\underline a^{-1},\quad
\beta=\underline u\,\underline b^{-1},\quad
\alpha_0=f^{-1}\alpha=f^{-1}\underline u\,\underline a^{-1},\quad
\beta_0=f^{-1}\deleted{\unsure{\alpha}}\inserted{\beta}=f^{-1}\underline u\,\underline b^{-1}.
\]
\mspage{453}

\scanpage{455}
\mspage{530}

On aura d'ailleurs
\[
\deleted{\unsure{\det\underline U=}}\det\underline u
=\det\underline U=\mu
\]
\[
\begin{aligned}
\det\alpha&=\underline{\det}(\underline u)\det(\underline a)^{-1}=\varepsilon'
&&\text{par (152')}\\
\det\beta&=\det(\underline u)\det(\underline b)^{-1}=1\\[3pt]
\text{et}\quad\det\alpha_0&=\det\alpha=\varepsilon'\\
\det\beta_0&=\det\beta=1
\end{aligned}
\]
on aura donc bien
\[
\begin{aligned}
\alpha,\alpha_0&\in\Gamma G^{\pm1}\overset{\text{déf}}=
\{g\mid(\det g)^2=1\}\\
\beta,\beta_0&\in\Gamma SG
\end{aligned}
\]
--- mais on a admis, bien sûr, que les \unsure{hom.} $\delta_\rho,\delta_\sigma$ de (135) sont \unsure{induits}, via les plongements (155), par le déterminant $\delta_G$.

Pour parvenir à la relation \unsure{ou} \unsure{d'équation}
\[
(157)\qquad
\underbrace{[\alpha,\rho]}_{\xi}
=\varepsilon\underbrace{[\beta,\sigma]}_{\eta}\varepsilon_1(\mu-1),
\]
il faut se donner des sections $\rho,\sigma$ de $G$ (\deleted{\unsure{ou}} plus précisément sans doute, \inserted{resp.\textsuperscript{t}} de $T_\rho$ et de $T_\sigma$ maximaux), et un unipotent strict $\varepsilon_1\in\Gamma SG$, pour que la relation (157) prenne un sens.
Remplaçons $\alpha,\beta$ par \unsure{des|les} valeurs $\underline u\,\underline a^{-1},\underline u\,\underline b^{-1}$, la relation devient (utilisant la commutativité de $T_\rho$, \deleted{\unsure{et}} qu'on n'avait pas vraiment fait \unsure{intervenir} jusqu'ici)
\[
(*)\qquad
[u,\rho]
=\varepsilon\cdot\underline u
\bigl(\underline b^{-1}\sigma(\underline b)^{\deleted{\unsure{-1}}}\bigr)^{\deleted{\unsure{-1}}}
\sigma(\underline u)^{-1}\varepsilon_1(\mu-1).
\]
\mspage{454}

\scanpage{456}

et résultera de l'ensemble des deux conditions
\[
(158)\quad
\left\{
\begin{aligned}
\text{a)}\quad&b^{-1}\sigma(b)=\varepsilon_\sigma(b)
&&\text{\uline{identiquement} si }b\in\Gamma N'_\sigma\\
&\text{i.e.}\quad b(\sigma)=\varepsilon_\sigma(b).\sigma
&&\text{identiquement}\\[4pt]
\text{b)}\quad&[u,\rho]=[u,\sigma]\varepsilon_1(\mu-1)
&&\text{identiquement en }u\in\Gamma T_0\\
&&&\text{pour }\mu=\delta_{T_0}(u)\overset{\mathrm{dfn}}=\det(u)\\
&\text{i.e.}\quad\rho(u)=\varepsilon_1(1-\mu)\sigma(u)
&&\text{identiquement en }u
\end{aligned}
\right.
\]
ce qui s'écrit aussi (en appliquant $\sigma^{-1}$ aux deux \unsure{mbs|mb})
\[
\sigma^{-1}\rho(u)=\varepsilon_0(1-\mu).u
\qquad\text{où}\quad\varepsilon_0=\sigma^{-1}(\varepsilon_1)
\]
et n'est guère concevable que si \uline{$T_0$ normalise $L_0$}, et si $\sigma^{-1}\rho\in\Gamma T_0.L_0$ \deleted{et} est de la forme
\[
(159)\quad
\left\{
\begin{aligned}
\text{a)}\quad&\widetilde{\sigma^{-1}\rho}=\widetilde{\varepsilon_0u_0}
\qquad u_0\in\Gamma T_0\\
\text{b)}\quad& T_0\text{ opère sur }L_0\text{ par la \unsure{coaction|action} \unsure{induite}}\\
&\hspace{8em}\text{par le déterminant}
\end{aligned}
\right.
\]
(Il y a quand $\unsure{\overline{\mathrm m}}$ encore de la marge, \unsure{parr|par} aux conditions draconiennes envisagées dans \uline{$\overline{\mathrm{VII}}$} !)

Mais ici l'équation (157) venait un peu comme cheveux sur soupe --- il doit y avoir des façons plus convaincantes et exigeantes de l'amener savoir \deleted{par} en relation avec des formules
\[
(160)\quad
\left\{
\begin{aligned}
\underline U(\rho)&=\operatorname{int}(\alpha_0).\rho
&&\bigl(=\alpha_0\rho(\alpha_0)^{-1}\rho\bigr)\\
\underline U(\deleted{\unsure{\rho}}\inserted{\sigma})&=
\varepsilon\operatorname{int}(\beta_0)\sigma
&&\bigl(=\beta_0\sigma(\beta_0)^{-1}\bigr)\\
\underline U(\varepsilon_0)&=\operatorname{int}(f^{-1})
\bigl(\varepsilon_0^{\deleted{\unsure{\mu}}}\unsure{(\mu)}\bigr)
\end{aligned}
\right.
\qquad
\begin{aligned}
\text{i.e.}\quad[\underline U,\rho]&=\alpha\rho(\alpha)^{-1}\\
[\underline U,\sigma]&=\varepsilon\beta\sigma(\beta)^{-1}
\end{aligned}
\]
\margin{attention au signe $\varepsilon$}

La troisième relation résulte aussitôt de $\underline U=f^{-1}\underline u$, et du fait que par hypothèse \inserted{(159 b)} $\underline u\in\Gamma T_0$ opère sur $L_0$ (\unsure{identiquement}) avec le multiplicateur $\mu=\det(\underline u)$.
\mspage{455}

\scanpage{457}
\mspage{531}

Les deux autres formules s'\unsure{explicitent} successivement, en écrivant $\alpha=\underline u\,\underline a^{-1}$, $\beta=\underline u\,\underline b^{-1}$, $\underline U=f^{-1}\underline u$
\[
\alpha_0=f^{-1}\alpha,\quad\beta_0=f^{-1}\beta
\]
\[
(161)\qquad
\left\{
\begin{aligned}
\underline u(\rho)&=\operatorname{int}(\alpha)(\rho)\\
\underline u(\sigma)&=\varepsilon\operatorname{int}(\beta)(\sigma)
\end{aligned}
\right.
\qquad\text{où }f\text{ ne figure plus}
\]
\[
\begin{aligned}
\text{i.e.}\qquad\underline u(\rho)&=\operatorname{int}(\underline u)(\rho)
&&\text{puisque }\operatorname{int}(\underline a^{-1})\rho=\rho\\
&&&\text{(comme }T_\rho\text{ comm.)}\\
&\text{OK !}\\
\underline u(\sigma)&=\deleted{\illegible{1}}\inserted{\varepsilon}\operatorname{int}(u)
\underbrace{\bigl(\underline b^{-1}(\sigma)\bigr)}_{\substack{=\varepsilon.\sigma\\\text{cf (158)(a)}}}
=\operatorname{int}(u)\sigma\quad\text{OK !}
\end{aligned}
\]

Pour prouver (160) on n'a utilisé que (158 a) et (159) (b). Essayons d'en déduire maintenant la ``relation des lacets'' (157) --- ça doit être un calcul formel, que j'ai déjà dû faire une douzaine de fois, dans des contextes qui se \unsure{différencient} au fur et à mesure ! On élimine $f$ en écrivant les relations (160) sous la forme équivalente (161) et
\[
(161\ \text{bis})\quad
\underline u(\varepsilon_0)=\varepsilon_0(\mu)
\qquad\text{i.e.}\quad
\underbrace{[\underline u,\varepsilon_0]}_{u\varepsilon_0(u)^{-1}}
=\varepsilon_0(\mu-1)\qquad\deleted{\unsure{\operatorname{int}}}
\]
\[
\deleted{\text{i.e.}\quad u\varepsilon_0(u)^{-1}}
\]
\[
\begin{array}{l|ll}
\begin{aligned}[t]
\text{or}\qquad&\\
[\underline u,\varepsilon_0]&=
u\underbrace{\varepsilon_0(u)^{-1}}_{\substack{\unsure{\text{i.e.}}\ \sigma^{-1}\rho(u)^{\deleted{\unsure{-1}}}}}\\
\text{i.e.}\quad\sigma([u,\varepsilon_0])&=\sigma(u)\rho(u)^{-1}
\end{aligned}
&
\begin{aligned}[t]
\text{or}\quad&\deleted{\rho(u)=\rho u\rho^{-1}=}\\
&u\rho u^{-1}\rho^{-1}=\xi\\
&\rho(u^{-1})=u^{-1}\xi\\
&\rho(u)=\xi^{-1}u\\
\unsure{\text{de}=|\text{de }\overline{\mathrm m}}\quad&\sigma(u)=\deleted{\unsure{\varepsilon}}\eta^{-1}u
\end{aligned}
&
\begin{aligned}[t]
\xi&\overset{\mathrm{dfn}}{=}\alpha\rho(\alpha)^{-1}\\[4em]
\eta&\overset{\mathrm{dfn}}{=}\varepsilon\beta\sigma(\beta)^{-1}
\end{aligned}
\end{array}
\]
\[
\text{d'où}\qquad
\sigma([u,\varepsilon_0])=
\deleted{\unsure{\xi^{-1}u\eta^{-1}u}}\eta^{-1}\xi
\]
et \unsure{(*)|($\alpha$)} devient, en appliquant $\unsure{\sigma}$
\[
\eta^{-1}\xi=\varepsilon_1(\mu-1)
\qquad\text{i.e.}\quad\xi=\eta\varepsilon_1(\mu-1)
\qquad\text{i.e. (157)}
\]
\hfill OK !
\mspage{456}

\scanpage{458}

En résumé :

\uline{Proposition} Supposons donnés, en plus des données de la section précédente $7^\circ$ (p. 526' ff \unsure{\}|)}, notamment (135)(136)), \deleted{des sections $\rho$ de $\unsure{T_{\rho,\sigma}|T_\rho,\sigma}$} ceci

\diagram{
\[
(162)\quad\left\{
\begin{array}{ll}
\text{a)}&\text{section }\rho\in\Gamma T_\rho,\quad\sigma\in\Gamma T_\sigma\quad\text{sections \uline{centrales}}\\[4pt]
\text{b)}&\text{Plongements}\quad
\begin{tikzcd}[ampersand replacement=\&,column sep=large]
T_\rho \arrow[r,hook] \& G
\end{tikzcd},\quad
\begin{tikzcd}[ampersand replacement=\&,column sep=large]
N'_\sigma \arrow[r] \& G
\end{tikzcd}\\[4pt]
\text{c)}&\text{tore}\quad
\begin{tikzcd}[ampersand replacement=\&,column sep=large]
T_0 \arrow[r,hook] \& G
\end{tikzcd}\quad\text{tel que }\delta_G:\quad
\begin{tikzcd}[ampersand replacement=\&,column sep=large]
\unsure{T_0} \arrow[r,"\sim"] \& G_m
\end{tikzcd}\\
&\text{i.e.}\quad G\simeq T_0.SG\quad\text{(semi-direct)}\\[4pt]
\text{d)}&\varepsilon_0\in\Gamma G\quad\text{unipotent \deleted{strict} régulier, définissant}\\
&\text{donc}\quad
\begin{tikzcd}[ampersand replacement=\&,column sep=large]
\unsure{\lambda} \arrow[r,mapsto] \& \varepsilon_0(\lambda)
\end{tikzcd}:\quad
\begin{tikzcd}[ampersand replacement=\&,column sep=large]
G_a \arrow[r,hook] \& SG
\end{tikzcd}
\end{array}\right.
\]
}

et supposons \deleted{(en plus de la \uline{commutativité} de $T_\rho,T_\sigma$}
\[
\begin{aligned}
\inserted{\text{a)}\quad\delta_G|T_\rho=\delta_\rho,\quad\delta_G|N'_\sigma=\delta_\sigma}\\
\deleted{\text{a)}\quad T_0\text{ \unsure{normalise} }\unsure{L_0}}
\end{aligned}
\]
\[
(163)\quad
\begin{aligned}
\text{b)}\quad&\operatorname{int}(\underline b)(\sigma)=\varepsilon_\sigma(\underline b)\sigma
\quad\text{dans }G\text{ identiquement}\\
&\hspace{12em}\text{pour }\underline b\in\Gamma\unsure{N'_\sigma|N^2_\sigma}\\[4pt]
\text{c)}\quad&T_0\text{ normalise }L_0,\text{ et on a}\\
&u(\varepsilon_0(\lambda))=
\varepsilon_0\bigl(\deleted{\unsure{\delta_G(u)}}\mu\lambda\bigr)
\quad\text{où }\mu=\delta_G(u)\\[4pt]
\text{d)}\quad&\widetilde{\unsure{\sigma|\rho}}^{-1}\widetilde\rho
\in\deleted{\unsure{\widetilde T_0}}\Gamma\widetilde\varepsilon_0\widetilde L_0
\quad\text{i.e. }\exists\text{loc.}^{\text{t}}\text{ section }\lambda_0\in\Gamma G_m\\
&\text{avec}\quad\lambda_0\widetilde\sigma^{-1}\widetilde\rho=\varepsilon_0u_0,
\quad\text{avec }u_0\in\Gamma T_0
\end{aligned}
\]
\margin{\unsure{relie} $L_0,T_0$}
\margin{$\begin{array}{l|}\text{\unsure{relie}}\quad L_0,T_0\\\rho,\sigma\end{array}$}

Sous ces conditions, on a pour tt système (152) satisfaisant les relations (152') (i.e. en fait pour tte \deleted{\unsure{syst}} section de $\mathcal M_{\rho,\sigma}$ !) un unique \unsure{syst.} (153), satisfaisant les relations (154) --- savoir $\underline U$ se met de façon unique sous la forme
\[
(164)\quad\underline U=f^{-1}\underline u,\qquad f\in\Gamma SG,\quad\underline u\in\Gamma T_0
\]
et on pose
\[
\begin{aligned}
(165)\quad\alpha&=\underline u\,\underline a^{-1},\quad
\beta=\underline u\,\underline b^{-1},\quad
\alpha_0=f^{-1}\unsure{\alpha|\alpha_\rho}=
\deleted{\unsure{\beta}}\inserted{f}^{-1}\underline u\,\underline a^{-1},\\
\beta_0&=f^{-1}\beta=f^{-1}\underline u\,\underline b^{-1}.
\end{aligned}
\]
\mspage{457}

\scanpage{459}
\mspage{532}

Tout ceci posé, on a alors les formules
(160), (161), (161 bis), et la formule ``du lacet'' (157).

Nous nous posons maintenant la question
sous quelles conditions la connaissance de
$(\alpha,\beta,f)$ (p.ex.) parmi les quantités (153),
permet de récupérer les quantités (152) ---
i.e. une section de $\mathcal M_{\rho,\sigma}$ --- et quelles conditions
il faut imposer à $(\alpha,\beta,f)$. Il
suffirait \inserted{(pour l'unicité)} de pouvoir construire $\underline u$, d'où
$\underline U,\underline a,\underline b$ par les trois premières formules
(154). \deleted{\unsure{Si, d'autre part on a}} Il faut
donc trouver $\underline a\in\Gamma T_\rho$ tel que $\alpha\underline a$ soit section
de $T_0$, --- indétermination sur $\underline a$ est de
$\Gamma T_0\cap T_\rho$, et $\underline b\in\Gamma N'_\sigma$ tel que $\beta\underline b$ soit section
de $T_0$ --- indétermination dans \deleted{$\unsure{\Gamma N'_\sigma\cap}$} $\Gamma T_0\cap N'_\sigma$.
L'indétermination sur $\underline u$ est la multiplication
à droite par une section de $T_0\cap T_\rho\cap N'_\sigma$, il
suffira donc qu'on ait
\[
\text{(164)}\qquad T_0\cap T_\rho\cap N'_\sigma=\{1\}.
\]
\margin{en fait, on aura $\unsure{\overline{\mathrm m}}$ $T_0\cap T_\rho=\{1\}$ dans les cas \unsure{envisagés} dans \uline{II}, \uline{III} et \unsure{suivants|suite}\ldots}
Supposant cette relation satisfaite, quelles
conditions mettre sur $(\alpha,\beta,f)$ pour provenir
d'un \unsure{syst.} (152) --- i.e. d'une section de $\mathcal M_{\rho,\sigma}$ ?

\mspage{458}

\scanpage{460}

On voit que \unsure{si} une section \deleted{$\unsure{\varphi}$} $\varphi_0$ de $\mathcal M_{\rho,\sigma}$, \deleted{est}
\deleted{\unsure{rempl.}} correspondant à $\alpha_0,\beta_0,f_0$, est remplacé
par $\varphi_1\overset{\unsure{\text{déf}}}{=}f^{-1}.\varphi_0$, correspondant à $\alpha_1,\beta_1,f_1$, alors
\[
\text{(165)}\qquad \alpha_1=\alpha_0,\quad\beta_1=\beta_0,\quad f_1=f_0f,
\]
car on voit que si $\varphi_0$ correspond à $(\underline U_0,\underline a_0,\underline b_0)$,
alors $\varphi_1$ à $(U_1,\underline a_1,\underline b_1)$, alors on a
\[
\text{(166)}\qquad U_1=f^{-1}U_0,\quad\underline a_1=\underline a_0,\quad\underline b_1=\underline b_0
\]
(car $\Theta_G\mid SG=$ l'inclusion canonique $SG\hookrightarrow G$,
et $\Theta_\rho$\inserted{$,\Theta_\sigma$} \deleted{est} triviaux sur $SG$). \deleted{Donc Cela
montre que} \inserted{Donc \deleted{\unsure{il}}} on aura
\[
u_1=u_0\qquad (\text{car }\det U_1=\det U_0)
\]
d'où (165). Cela montre que les conditions
cherchées sur $(\alpha,\beta,f)$ ne dépendent que de
$\alpha,\beta$, pas de $f$. On voudrait qu'elles
s'expriment par la validité de la formule des
lacets\inserted{(157)}, pour un paramètre \inserted{$\varepsilon,$} $\mu$ non précisés $(\varepsilon^2=1)$
(qui \unsure{sera} alors uniquement déterminé --- et \unsure{permettra}
la \unsure{définition sec.} de \unsure{notion locale|action locale}\ldots).

\margin{à vrai dire, cette condition n'est pas \uline{suffisante} $\overline{\mathrm m}$ dans les cas envisagés --- on trouvera \unsure{donc|dans} un $\underline u\in B'_0$, et il faut exiger que $\underline u\in T_0$ (déjà pour $(\alpha,\beta)$ soit ``normalisé'') i.e. $B=0$, où $B=ac+bd+bc$, si $\underline u=\left(\begin{smallmatrix}a&b\\c&d\end{smallmatrix}\right)$ (cas modulaire) cf prop. p.~412'\ldots\ \unsure{Pour l'avoir} $\unsure{R}$ ; \unsure{lui-même étant} \unsure{condition supplém.}, \unsure{il|qui} \unsure{faut remplacer|faut envisager} $\mathcal M_{\rho,\sigma}$ \unsure{par le produit} $\tfrac12$ \unsure{direct par} $L_0$, \unsure{sur lequel on fait} $\unsure{f_\Delta}=\unsure{0}$, \unsure{soit} $S\mathcal M_{\rho,\sigma}$.}

La suffisance de cette condition \unsure{fera}
l'objet des réflexions des (II) (III),
\unsure{sous} des conditions (plus draconiennes que celles
\unsure{de la} prop. précédente (p.~531') que nous
avons imposées au triplet $(\varepsilon_0,\rho,\sigma)$. Il

\mspage{459}

\scanpage{461}
\mspage{533}

est possible que les conditions auxquelles nous
\inserted{\unsure{y}} sommes parvenus étaient plus restrictives
que nécessaires, et que je sois obligé
d'y revenir par la suite, dans un contexte
plus général que celui des formes de
$\mathrm{Gl}(2)$, $\mathrm{GP}(1)$, $\mathrm{Sl}(2)$, mais pour le moment
l'étude détaillée faite en nous inspirant
des hypothèses dégagées dans \deleted{\unsure{les}} II, III nous
aura servi de déblayage, pour nous éclairer
les idées sur un cas particulier
important.

\uline{Rem.} Rappelons que nous avons vu qu'en
fait $\mu$ est déjà défini en termes de $\alpha$ seul,
et $\overline{\mathrm m}$ en termes de $\alpha\rho(\alpha)^{-1}=\xi$ \unsure{;} \deleted{\unsure{bien|donc}}
Mais il n'est déterminé en termes de $\beta$
--- en fait, de $\beta\sigma(\beta)^{-1}$ --- que modulo
la connaissance de $\varepsilon$ (en pratique, modulo
signe\ldots). Mais si l'on impose (p.ex.) $\varepsilon=1$,
la donnée de $\alpha$ ne détermine $\beta$ que
modulo \deleted{$\unsure{\beta}$} $\beta\longrightarrow\beta T_\sigma$\ldots

\mspage{460}

\scanpage{462}

La structure de $\mathcal M_{\rho,\sigma}$ étant bien comprise,
revenons sur le groupe $S\mathcal M_{\rho,\sigma}$\inserted{(cf p.~495')}, considéré comme
produit semi-direct de $\mathcal M_{\rho,\sigma}$ par le sous-groupe
$L_0$ de $SG$, sur lequel en fait opère $\mathcal M_{\rho,\sigma}$
via $\chi$. Désignant par

\begin{samepage}
\noindent (165)\quad
\diagram{\begin{tikzcd}[ampersand replacement=\&,column sep=large]
\deleted{\unsure{\varepsilon_0}:\unsure{U}}\;\inserted{\unsure{\mathcal E|\varepsilon}_{L_0}:G_a} \arrow[r,"\sim"] \&
\operatorname{Ker}\bigl(S\mathcal M_{\rho,\sigma}\longrightarrow\mathcal M_{\rho,\sigma}\bigr)
\end{tikzcd}}
\end{samepage}

\inserted{correspondant à $\varepsilon_0$,}
on aura par \deleted{\unsure{définition}} construction
\[
\text{(166)}\qquad
U\bigl(\unsure{\mathcal E|\varepsilon}_{L_0}(\lambda)\bigr)
=\unsure{\mathcal E|\varepsilon}_{L_0}(\mu\lambda)
\qquad\text{où }\mu=\chi(U).
\]

\deleted{Le groupe $S$}

\mspage{461}

\scanpage{463}
\mspage{534}

\deleted{\unsure{XI}} \textcircled{$9^\circ$} \uline{Les sous-groupes}
$\mathcal M_{\rho,\sigma}(0)$,
\deleted{$\mathcal M_{\rho,\sigma}(\tfrac12),\mathcal M_{\rho,\sigma}(-\unsure{\overline J|\overline\jmath})$.}
\inserted{$\mathcal M_{\rho,\sigma}[0]$, et le groupe}
$S\mathcal M_{\rho,\sigma}$, l'isom.
\diagram{\begin{tikzcd}[ampersand replacement=\&,column sep=large]
S\mathcal M_{\rho,\sigma} \arrow[r,"\sim"] \& \Sigma^+_{\rho,\sigma}
\end{tikzcd}}

Le groupe $\mathcal M_{\rho,\sigma}$ est extension de $\mathcal T_{\rho,\sigma}$ par
$\unsure{\widetilde G|\widetilde{\mathcal C}}^+_{\rho,\sigma}\simeq SG$, et les trois sous-groupes que nous allons
expliciter seront des \deleted{sections}\inserted{\unsure{ss-groupes|sous-groupes}}, \deleted{partielle} avec le
grain de sel que pour le troisième, ce sera une
section partielle seulement, au dessus du
sous-groupe \inserted{$\mathcal T'_{\rho,\sigma}$} d'indice $2$ de $\mathcal T_{\rho,\sigma}$, noyau de
$\varepsilon_3$.

$[\![$Déterminons de façon générale les sections
de $\mathcal M_{\rho,\sigma}$ sur $\mathcal T_{\rho,\sigma}$. Comme on a

\begin{samepage}
\noindent (1)\quad
\diagram{\begin{tikzcd}[ampersand replacement=\&,column sep=large,row sep=small]
R_{\rho,\sigma} \arrow[r] \arrow[d,phantom,"{\unsure{\wr|\simeq}}" description] \& \mathcal T_{\rho,\sigma} \quad \text{épi de noyau }\mu_2\\
S'T_\rho\times SN'_\sigma\times G_m \&
\end{tikzcd}}
\end{samepage}

\margin{finalement ça servira pas à grand chose\ldots}

il revient au $\overline{\mathrm m}$ de se donner une section,
ou : un hom.

\begin{samepage}
\noindent (2)\quad
\diagram{\begin{tikzcd}[ampersand replacement=\&,column sep=large]
R_{\rho,\sigma} \arrow[r,"{\deleted{\unsure{\Xi|\zeta}}\inserted{\psi}}"] \& \mathcal M_{\rho,\sigma}
\end{tikzcd}}
\end{samepage}

qui soit \deleted{trivial}
\[
\text{(3)}\qquad
\left\{\begin{aligned}
\text{a)}\quad &\text{trivial sur }\mu_2\\
\text{b)}\quad &\deleted{\unsure{\Xi|\Sigma}(x)\unsure{\equiv|=}x\ \bigl(\unsure{\widetilde G|\widetilde{\mathcal C}}^+_{\rho,\sigma}\bigr)\quad\text{identiquement}.}\quad\inserted{\text{``remonte'' (1)}}
\end{aligned}\right.
\]

Les hom. $\unsure{\Xi|\Sigma}$ qui satisfont à (3) correspondent
aux scindages de l'extension \uline{triviale} de
$R_{\rho,\sigma}$ par $SG$ (extension \uline{produit}), donc
simplement aux \uline{homomorphismes}

\begin{samepage}
\noindent (4)\quad
\diagram{\begin{tikzcd}[ampersand replacement=\&,column sep=large]
R_{\rho,\sigma} \arrow[r,"{\deleted{\unsure{\Xi|\zeta}_{\unsure{0}}}\inserted{\psi^0}}"] \& SG
\end{tikzcd}}
\end{samepage},

\mspage{462}

\scanpage{464}

en associant à $\psi^0\deleted{\unsure{\xi_0}}$ le hom.
\[
(5)\qquad x\longmapsto \deleted{\unsure{\xi}}\psi(x)
=x\deleted{\unsure{\xi}}\psi^0(x)
=\deleted{\unsure{\xi(x)x}}\psi^0(x).x
\]
Donc la condition a prend la forme
\[
(6)\qquad \deleted{\unsure{\xi}}\psi^0_{\unsure{0}}(x)=x
\quad\text{sur }\mu_2=R_{\rho,\sigma}\cap SG.
\]

\margin{\deleted{\unsure{$\rho$}}}
Les $\psi,\psi^0$ correspondants aux \deleted{\unsure{« trois »}} sections
$\mathcal M_{\rho,\sigma}^{(0)}$, $\mathcal M_{\rho,\sigma}^{(\frac12)}$,
$\mathcal M_{\rho,\sigma}^{(-\overline J)}$ seront notés resp.\textsuperscript{t}
\[
(\psi_{\unsure{0}},\psi^0_0),\ (\psi_{1/2},\psi^0_{1/2}),\
(\psi_{-\overline J},\psi^0_{-\overline J}).\qquad\big]\!\big]
\]

1°) \uline{Cas de} $\mathcal M_{\rho,\sigma}^{(0)}$. Avec les notations
de la p. 529, il \unsure{correspondra} \deleted{\unsure{$\underline U$ et les}}
\inserted{\unsure{$\underline u$}} de $\mathcal M_{\rho,\sigma}$ satisfaisant à l'une des conditions
\uline{équivalentes}
\[
(7)\qquad f=1,\quad\inserted{\underline U=\underline u,}\quad
\underline U\in\Gamma T_0,\quad \alpha_0=\alpha,\quad\beta_0=\beta
\]
donc on a
\[
(8)\qquad \mathcal M_{\rho,\sigma}^{(0)}=\Theta^{-1}(T_0).
\]
\inserted{(en termes de la description de $\mathcal M_{\rho,\sigma}$ de p. 526' ff)}

Vérifions que c'est un \unsure{sous-groupe-section}.
On a
\[
\overbrace{\mathcal M_{\rho,\sigma}^{(0)}\cap\widetilde{\unsure{G|\mathcal C}}_{\rho,\sigma}^{+}}
^{\inserted{(\Theta|SG)^{-1}(T_0)}}=\{1\}
\]
\margin{\unsure{et}}
car $\Theta|(\widetilde{\unsure{G|\mathcal C}}_{\rho,\sigma}^{+}=SG)$ est l'inclusion
canonique $SG\hookrightarrow G$, or $SG\cap T_0=\{1\}$.
Reste à vérifier que
\diagram{\begin{tikzcd}[ampersand replacement=\&]
\mathcal M_{\rho,\sigma}^{(0)} \arrow[r] \& \mathcal T_{\rho,\sigma}
\end{tikzcd}}
(qui est mono par \unsure{$(a)$|$(\alpha)$}) est épi. Mais en

\mspage{463}

\scanpage{465}
\mspage{535}

termes de la \unsure{réalisation} can.
\[
(9)\qquad\mathcal M_{\rho,\sigma}\simeq
\underbrace{R_{\rho,\sigma}\times SG}_{\widetilde{\mathcal M}_{\rho,\sigma}}
/\underset{\text{diagonal}}{\mu_2}
\]
et de l'hom
\margin{\diagram{\begin{tikzcd}[ampersand replacement=\&]
\kappa:R_{\rho,\sigma}\arrow[r]\&\mathbf G_m
\end{tikzcd}}
l'hom canonique \unsure{permettant} d'exprimer $\Theta|R_{\rho,\sigma}$}
\diagram{\begin{tikzcd}[ampersand replacement=\&]
(10)\qquad i_{T_0}:\mathbf G_m\arrow[r,"\sim"]\&T_0\arrow[r,hook]\&G
\end{tikzcd}}
\unsure{isom|hom} inverse de l'isom
\diagram{\begin{tikzcd}[ampersand replacement=\&]
\delta_{T_0}=\det|T_0:\quad T_0\arrow[r,"\sim"]\&\mathbf G_m,
\end{tikzcd}}
on a, pour un élément
\diagram{\begin{tikzcd}[ampersand replacement=\&,column sep=1em,row sep=.2em]
(\underline r,\&U_0)\\
\inserted{R_{\rho,\sigma}}\arrow[u]\&\inserted{SG}\arrow[u]
\end{tikzcd}}
de $\Gamma\widetilde{\mathcal M}_{\rho,\sigma}$
\[
\Theta(\underline r,U_0)=\kappa(\underline r)U_0,\quad
\text{déterminant }\kappa(r)^2
\]
dire que c'est $\unsure{\in U}$ signifie donc l'on a
\[
\kappa(\underline r)U_0=i_{T_0}(\kappa(\underline r)^2)
\quad\text{i.e. }\deleted{\unsure{U_0=}}
\]
\[
(11)\qquad U_0=\kappa(\underline r)^{-1}i_{T_0}(\kappa(\underline r)^2)
\]
donc l'image inverse $\widetilde{\mathcal M}_{\rho,\sigma}^{(0)}$ de
$\mathcal M_{\rho,\sigma}^{(0)}$ dans $\widetilde{\mathcal M}_{\rho,\sigma}$ est formée des
\[
(12)\qquad
\bigl(\underline r,\kappa(\underline r)^{-1}i_{T_0}(\kappa(\underline r))^2\bigr)
\in\Gamma\widetilde{\mathcal M}_{\rho,\sigma}
\]
en d'autres termes cela correspond à un \unsure{groupe} section, avec
\[
(13)\qquad\psi^0_0(\underline r)=\kappa(\underline r)^{-1}i_{T_0}(\kappa(\underline r)^2).
\]
Comme il se doit, on a bien là un \uline{hom}
\diagram{\begin{tikzcd}[ampersand replacement=\&]
\psi^0_0:R_{\rho,\sigma}\arrow[r]\&SG
\end{tikzcd}}
et sa valeur sur $\mu_2$ est l'identité,
car $\kappa(\underline r)^2=1$ et on a
\[
\psi^0_0(r)=\kappa(\underline r)^{-1}=
\begin{gathered}\kappa(\underline r)\\\Vert\\\underline r\end{gathered}
\]

\uline{Remq}\quad\deleted{Avant de \unsure{qui}}
Pour définir cette section, il suffisait

\mspage{464}

\scanpage{466}

(au pt de vue du pur calcul) de se donner un hom quelconque, noté
$\unsure{i'_{T_0}|i_{T_0}}$, de $\mathbf G_m$ dans $SG$, et de poser (13).

Avant de quitter $\mathcal M_{\rho,\sigma}^{(0)}$, je vais introduire le groupe
un peu plus grand (de dim 4, au lieu de 3)
\[
(14)\qquad\mathcal M_{\rho,\sigma}^{[0]}
\overset{\mathrm{dfn}}{=}\Theta^{-1}(B'_0)
\]
où, on le rappelle, $B'_0$ est défini comme le produit $\frac12$ direct
\[
(15)\qquad B'_0=T_0.L_0
\]
Cette fois-ci, $U=U_0\kappa(r)\inserted{\in\Gamma B'_0}$ n'est pas déterminé
par la seule connaissance de son déterminant\inserted{\unsure{$2$}}
$\underline\kappa(r)^2$, qui ne le détermine que modulo $L_0$.
Cela correspond au fait que
$\underset{\inserted{=B'_0\cap SG}}{
\mathcal M_{\rho,\sigma}^{[0]}\cap
\overbrace{\widetilde{\unsure{G|\mathcal C}}_{\rho,\sigma}^{+}}^{\inserted{SG}}}$
n'est plus réduit à $(1)$, mais à $L_0$, qui apparaît donc comme
un sous-groupe invariant de $\mathcal M_{\rho,\sigma}^{[0]}$, et on \unsure{aura|a} évidemment
(compte tenu que $\mathcal M_{\rho,\sigma}^{(0)}\subset\mathcal M_{\rho,\sigma}^{[0]}$ et
\diagram{\begin{tikzcd}[ampersand replacement=\&]
\mathcal M_{\rho,\sigma}^{(0)}\arrow[r,"\sim"]\&\mathcal T_{\rho,\sigma}
\end{tikzcd}}) \unsure{par|que}
\[
(16)\qquad\begin{aligned}
\mathcal M_{\rho,\sigma}^{[0]}&=\mathcal M_{\rho,\sigma}^{(0)}.L_0\\
&\qquad\text{produit }\frac12\text{ direct}
\end{aligned}
\]
C'est encore une construction qui a un sens dès que l'on a un $T_0$ comme dessus,
\deleted{\illegible{2}} et un sous-groupe \deleted{\uline{\unsure{affine}}} $L_0$ de $SG$,

\mspage{465}

\scanpage{467}
\mspage{536}

\unsure{invarié} par $T_0$. \deleted{\unsure{On}}

L'opération de $\mathcal M_{\rho,\sigma}^{[0]}$ sur $SG$ se fait
par l'intermédiaire de $\Theta$, depuis peu regardé comme un épi
\diagram{\begin{tikzcd}[ampersand replacement=\&]
\mathcal M_{\rho,\sigma}^{[0]}
\arrow[r,"\Theta|\mathcal M_{\rho,\sigma}^{[0]}"]\&B'_0
\end{tikzcd}};
cette opération étant précisée, on peut introduire
\margin{pour une définition plus \unsure{transparente}, cf p. 538', 539}
\[
(17)\qquad\deleted{\unsure{\mathcal S}}
S\mathcal M_{\rho,\sigma}\overset{\text{déf}}{=}
\begin{gathered}
\mathcal M_{\rho,\sigma}^{[0]}.SG\\
\text{produit }\frac12\text{ direct}
\end{gathered}
\]
en lequel $SG$ est invariant, et $\mathcal M_{\rho,\sigma}^{[0]}$ s'identifie
à un sous-groupe de $S\mathcal M_{\rho,\sigma}$. Comme
$\mathcal M_{\rho,\sigma}^{(0)}\subset\mathcal M_{\rho,\sigma}^{[0]}$,
on a une inclusion can.
\[
(18)\qquad
\underbrace{\mathcal M_{\rho,\sigma}^{(0)}:SG}_{\unsure{\simeq\mathcal M_{\rho,\sigma}}}
\subset
\underbrace{\mathcal M_{\rho,\sigma}^{[0]}.SG}_{=S\mathcal M_{\rho,\sigma}}
\]
dont l'image est \deleted{\unsure{notée} $S\mathcal M_{\rho,\sigma}^{(0)}$}
identifiée à $\mathcal M_{\rho,\sigma}$.
\deleted{Ceci dit, le sous-groupe $L_0$ de $\mathcal M_{\rho,\sigma}^{[0]}$
cf (16) \unsure{opère} trivial\textsuperscript{t} sur $SG$,
i.e. $\Theta\inserted{|\mathcal M_{\rho,\sigma}^{[0]}}$,
qui est trivial en effet sur $L_0$}

\deleted{\unsure{Revenons|Reprenons}} je commence à \unsure{nouveau} à m'y perdre, je vais me raccrocher
--- pour une compréhension de la signification "géom." de
$S\mathcal M_{\rho,\sigma}$, $\mathcal M_{\rho,\sigma}^{[0]}$, aux notations
de la p. 529, qui mérite d'être \unsure{explicité|explicitée} par un

\mspage{466}

\scanpage{468}

\uline{Proposition}\quad Considérons l'hom
\diagram{\begin{tikzcd}[ampersand replacement=\&,column sep=large,row sep=large]
(19)\quad\mathcal M_{\rho,\sigma}\arrow[r]
\&G\times T_\rho\times N'_\sigma\times\mathbf G_m\times
\unsure{\mu_2^2|\mu_2}
\deleted{\unsure{(\mathbf G_m^3\times\mathbf G_m\times\mu_2)}}\\
x\arrow[r,mapsto]
\&(\Theta_G(x),\Theta_\rho(x),\Theta_\sigma(x),\chi(x),\varepsilon_2(x),\varepsilon_3(x))
\end{tikzcd}}
C'est un monomorphisme, dont l'image est
formée des \deleted{triplets}\inserted{\unsure{syst.}}
$(\underline U,\underline a,\underline b,\inserted{\mu,\varepsilon,\varepsilon'})$
satisfaisant à
\[
\deleted{\unsure{\mu=|\text{pour}}\,
\unsure{\det_G|\delta_G}(\underline U)
=\delta_\rho(\underline b)=\varepsilon\det}
\]
\[
(20)\qquad\left\{\begin{aligned}
\mu&=\delta_G(\underline U)=\delta_\sigma(\underline b)
=\varepsilon'\delta_\rho(\underline a)\\
\varepsilon&=\varepsilon_\sigma(\underline b)\\[1ex]
\end{aligned}\right.
\]

\uline{Corollaire}. Le morphisme.
\diagram{\begin{tikzcd}[ampersand replacement=\&,column sep=large,row sep=large]
{}\arrow[d,no head,"(21)"']\&\mathcal M_{\rho,\sigma}\arrow[r]
\&G\times T_\rho\times N'_\sigma\\
{}\&x\arrow[r,mapsto]\&\Theta_G(x),\Theta_\rho(x),\Theta_\sigma(x)
\end{tikzcd}}
est un monomorphisme, l'image est formée des triplets
$(\underline U,\underline a,\underline b)$ satisfaisant à
\margin{\uline{NB} on a
$\left\{\begin{aligned}
\varepsilon&=\varepsilon_\sigma(\underline b)\\
\varepsilon'&=\delta_\rho(\underline a)/\delta_\sigma(\underline b)\\
\mu&=\delta_\sigma(\underline b)=\delta_G(\underline U)
\end{aligned}\right.$}
\[
(22)\qquad\left\{\begin{aligned}
\delta_G(\underline U)&=\delta_\sigma(\underline b),\\
[\delta_\rho(\underline a)/\delta_\sigma(\underline b)]^2&=1.
\end{aligned}\right.
\]

\uline{Dém. de la proposition.}\quad On a
\deleted{$\unsure{\operatorname{Ker}(\Theta_G)|\mu_G}$}
\[
\operatorname{Ker}(\Theta_G,\deleted{\unsure{\chi|\xi_2}},\varepsilon_2,\varepsilon_3)
=SR_{\rho,\sigma}^{\unsure{\circ|0}}\simeq ST_\rho\times ST_\sigma
\]
et $\Theta_\rho,\Theta_\sigma$ \deleted{\unsure{donc|des}}
sur $SR_{\rho,\sigma}^{\unsure{\circ|0}}$ s'identifient aux deux projections,
donc \deleted{$\unsure{\Theta}$}
\[
\operatorname{Ker}\bigl((\Theta_\rho,\Theta_\sigma)|SR_{\rho,\sigma}^{\unsure{\circ|0}}\bigr)
=\{1\},\quad\text{donc}
\]
\[
\operatorname{Ker}(\Theta_G,\chi,\varepsilon_2,\varepsilon_3,\Theta_\rho,\Theta_\sigma)
=(1),
\]

\mspage{467}

\scanpage{469}
\mspage{537}

donc mono. Pour l'épi, on est ramené à la forme équivalente du Corollaire. --- OPS qu'on a un $\lambda$ tel que
\[
\lambda^2=\mu\overset{\mathrm{dfn}}{=}\det\underline U
\]
on veut écrire
\[
\underline U=\Theta_{\widetilde{\mathcal M}}
\left(\underset{\underset{R_{\rho\sigma}}{\rotatebox{90}{$\in$}}}{r},
\underset{\underset{SG}{\rotatebox{90}{$\in$}}}{U_0}\right)
=\kappa(r)U_0
\]
\margin{NB $\Theta_{\widetilde{\mathcal M}}$ l'hom can.
\diagram{\begin{tikzcd}[ampersand replacement=\&]
\widetilde{\mathcal M}_{\rho,\sigma}\overset{\mathrm{dfn}}{=}R_{\rho,\sigma}\times SG \arrow[r] \& G
\end{tikzcd}}
que j'\unsure{avais|écrivais|écris} $\Theta_G$.}
avec $U_0\in\Gamma(SG)$, i.e. $\kappa(r)^2=\det\underline U$ i.e. $\kappa(r)^2=\mu$,
on prendra $r$ avec \deleted{$\unsure{\kappa}$}
\[
\kappa(r)=\lambda,
\]
\deleted{\unsure{alors}} on pose donc
\[
\underline U_0=\lambda^{-1}U,\quad
\underline a_0=\lambda^{-1}\underline a,\quad
\underline b_0=\lambda^{-1}b
\]
\deleted{\unsure{et on est}} \inserted{et on est} ramené à voir que $\underline U_0,\underline a_0,\underline b_0$ sont dans l'image de $\mathcal M_{\rho,\sigma}$ (en \unsure{étant|ôtant} des ``\unsure{corrigés}'' $(\underline U,\underline a,\underline b)$ par l'image \deleted{\unsure{dans le}} $G\times T_\rho\times\unsure{T_\sigma}$ de la section $\lambda$ de $G_m\hookrightarrow R_{\rho,\sigma}$). On est ramené à voir que $(\underline a_0,\underline b_0)$ est dans l'image de
\diagram{\begin{tikzcd}[ampersand replacement=\&]
R_{\rho,\sigma}\arrow[r] \& \deleted{\unsure{T}}T_\rho\times N'_\sigma.
\end{tikzcd}}
Mais maintenant $\underline a_0$ est dans $S'T_\rho$, $\underline b_0$ dans $SN'_\sigma$, et justement $R_{\rho,\sigma}\simeq S'T_\rho\times SN'_\sigma$, \unsure{c'est gagné}\ldots

\uline{Remarque}. On ne peut remplacer $G$ par $PG$,
i.e.
\diagram{\begin{tikzcd}[ampersand replacement=\&]
\mathcal M_{\rho,\sigma}\arrow[r] \& PG\times T_\rho\times N'_\sigma
\end{tikzcd}}
(où on peut rajouter encore $G_m\times\mu_2\times\mu_2$ si on veut!) n'est \uline{pas} mono --- son noyau est justement le sous-groupe $\mu_2$ de $\widetilde{\unsure{G|\mathcal C}}_{\rho,\sigma}^{+}=SG$.

\mspage{468}

\scanpage{470}

Considérons, maintenant que les quantités $\underline U,\underline a,\underline b$ $(\inserted{et}\ \mu,\varepsilon,\varepsilon')$ de p.~529 (152) sont comprises, les quantités (153) du loc. cit $\underline u,f,\alpha,\beta,\alpha_0,\beta_0$ ($g,h\ \xi,\eta$ n'y figurent, en plus\ldots) Elles sont déterminées en termes des $\underline U,\underline a,\underline b$, dès qu'on connaît soit $\underline u$, soit $f$, par les formules (154), $\underline u$ et $f$ étant reliés par la seule relation
\[
\underline U=f^{-1}\underline u \tag{23}
\]
avec en plus sur $f$ la condition
\[
f\in\Gamma\underbrace{SG}_{\widetilde{\unsure{G|\mathcal C}}_{\rho,\sigma}^{+}}^{\unsure{\sim}},\tag{24}
\]
plus une condition sur $u$ qui, dans loc. cit, s'écrivait
\[
\underline u\in\Gamma T_0,\tag{25}
\]
et était suggérée par les calculs de \S46 $\overline{\underline{\mathrm V}}$ (repris ensuite dans \S48 $\overline{\underline{\mathrm{II}}}$, $\overline{\underline{\mathrm{III}}}$ etc). Cette condition assure l'\uline{unicité} de $f,\underline u$ \unsure{satisfaisant|satisfaisants} à (23) --- c'est là, du $\mathrm p^{\mathrm t}$ de vue du calcul, sa principale vertu! --- mais dans certains contextes il est \unsure{préférable} de la remplacer par la condition plus \unsure{large}
\[
\underline u\in\Gamma B'_0
\qquad\text{où}\quad
\begin{gathered}
B'_0=T_0.L_0.\\
\text{(produit }\tfrac12\text{ direct)}
\end{gathered}\tag{26}
\]

\mspage{469}

\scanpage{471}
\mspage{538}

On n'a plus \deleted{de \unsure{résultat|résultats}} d'\uline{unicité} pour l'écriture de $\underline U$ sous forme (23) --- par contre on a un théorème d'\uline{existence} sous forme simple :

Si on \uuline{part} de $(\alpha,\beta,f)$\deleted{$\unsure{\in}$}, sections de $G$ telles que
\[
(\det\alpha)^2=\det\beta=\det f=1,\tag{27}
\]
pour que \deleted{l'on puisse écrire} \inserted{ce triplet} et par \unsure{conséquent|conséquence} les quantités qui s'en déduisent : \deleted{$\alpha_0,\beta_0,$}
\deleted{$\xi,\eta,g,h$}
\[
\left\{
\begin{aligned}
\alpha_0&=f^{-1}\alpha,&\quad \beta_0&=f^{-1}\unsure{\varepsilon}\beta\\
\xi&=\alpha\rho(\alpha)^{-1},&\quad \eta'&=\beta\sigma(\beta)^{-1}\\
g&=\alpha_0\rho(\alpha_0)^{-1},&\quad h&=\beta_0\sigma(\beta_0)^{-1}
\end{aligned}
\right.\tag{28}
\]
proviennent d'un système $(x,\deleted{\unsure{f}}\inserted{\underline u})$,
avec $x\in\Gamma\mathcal M_{\rho,\sigma}$,
$\deleted{\unsure{f}}\inserted{\underline u}\in\Gamma\deleted{\unsure{SG|T_0}}\inserted{B'_0}$ tel que
\[
f\overset{\mathrm{dfn}}{=}\underline u\,\underline U^{-1}
\deleted{\unsure{\underline u|f}}
\in\Gamma\deleted{\unsure{B'_0}}\inserted{SG},
\quad\text{i.e.}\quad\det\underline u=\det\underline U,
\tag{29}
\]
par les formules
\[
\alpha=\underline u\,\underline a^{-1},\quad
\beta=\underline u\,\underline b^{-1},\quad
f=\underline u\,\underline U^{-1},\tag{30}
\]
il f. et s. que l'on ait la relation (qui ne fait pas intervenir $f$, seulement $\alpha,\beta$) :
\[
\alpha\rho(\alpha)^{-1}=\varepsilon\beta\sigma(\beta)^{-1}\varepsilon_1(\mu-1)
\qquad
\begin{gathered}
\text{pour : }\varepsilon\in\Gamma\mu_2\\
\mu\in\Gamma G_m\\
\text{convenables.}
\end{gathered}\tag{31}
\]

\uline{(NB} $\unsure{\varepsilon|\varepsilon_1},\mu$ sont uniquement déterminés (\unsure{donc} en termes de $\alpha\ \beta$).
\inserted{Alors le couple $(x,\underline u)$ est unique. Pour $f=1$, (i.e. $\underline u=\underline U$),}
On peut dire \deleted{que} aussi qu'on a une correspondance $1$-$1$ entre les \inserted{l'ens. \deleted{\illegible{1}} sections \unsure{du|des}} \unsure{sous-groupes|sous-groupe} $\mathcal M_{\rho,\sigma}[0]$ (pas $\mathcal M_{\rho,\sigma}(0)$!) de $\mathcal M_{\rho,\sigma}$, et l'ens.

\mspage{470}

\scanpage{472}

des couples $\alpha,\beta$ satisfaisant aux deux relations
\[
(\det\alpha)^2=\det\beta=1\qquad\text{de (27)}
\]
\uline{et} la relation des lacets. Pour cela, l'on n'a \deleted{\unsure{d'ailleurs}} pas d'utilité pour le groupe $S\mathcal M_{\rho,\sigma}$.

Mais revenant \unsure{au cas} $f$ quelconque, pour interpréter le \unsure{faisceau} \unsure{algébrique} sur les triples $(\alpha,\beta,f)$, \inserted{satisfaisant (27)(31),} \inserted{avec la structure de} groupe étrange dessus \inserted{\unsure{via} relation de cocycles,} (sur lequel il me faudra revenir en termes \inserted{plus} \unsure{géométriques}\ldots) et parce que les trois quantités $\alpha,\beta,f$ se mélangent intimement --- il y a lieu d'introduire $S\mathcal M_{\rho,\sigma}$ --- qui sera isomorphe \unsure{au|à} schéma (en groupes) des solutions communes de (27)(31).

De ce point de vue, $S\mathcal M_{\rho,\sigma}$ se déduit de $\unsure{\rho}$ comme \deleted{\unsure{décrit par triples}} \inserted{schéma}
\deleted{$(x,\underline u,f)$ \unsure{de}.}
\[
\begin{aligned}
S\mathcal M_{\rho,\sigma}&=\text{schéma des triples }(x,\underline u,f)\\
&\text{avec }\left\{
\begin{aligned}
x\ \deleted{\unsure{\in}}&\text{ dans }\mathcal M_{\rho,\sigma},\\
\deleted{\unsure{f}}\inserted{\underline u}&\text{ dans }B'_0\\
f&\text{ dans }SG
\end{aligned}
\right.
\qquad\left\{
\begin{aligned}
\underline U&=f^{-1}\underline u\\
\Vert\,\text{déf}&\\
\Theta_G(x)&
\end{aligned}\right.
\end{aligned}\tag{32}
\]

Donc
\diagram{\begin{tikzcd}[ampersand replacement=\&,column sep=large,row sep=small]
{(33)\quad S\mathcal M_{\rho,\sigma}}
\arrow[r,hook]
\arrow[d,phantom,"{\unsure{\rotatebox{90}{$\simeq$}}}"]
\& {\mathcal M_{\rho,\sigma}\times B'_0\quad\unsure{\leftarrow}}\\
{\left\{x,\underline u\ \middle|\ \unsure{\chi|\gamma}(x)=
\underbrace{\unsure{\delta_G}(\underline u)}_{\begin{gathered}\Vert\,\mathrm{dfn}\\\det\underline u\end{gathered}}\right\}} \&
\end{tikzcd}}
ou encore, plus simplement

\mspage{471}
